\section{Methods}
\label{sec:methods}

\begin{figure}[htbp]
\centering
\includegraphics[width=\textwidth]{fig01_schematic}
\caption{(a) Cross-section of a top-fired box reformer: burners in the roof between rows of catalyst tubes, flames releasing heat
in the upper part of the firebox, flue gas flowing downward co-currently with the process gas and leaving through tunnels at the
floor. (b) Module chain of the digital twin: operating inputs are mapped through the furnace and tube models, which exchange the
outer-wall temperature profile and the heat flux, to the creep-life module; surrogates and uncertainty quantification wrap the chain.}
\label{fig:schematic}
\end{figure}

The twin is implemented as a Python package (\texttt{rdt}) whose modules follow the chain of Fig.~\ref{fig:schematic}b. All
physical data (kinetic parameters, plant operating points, creep curves) are read from annotated YAML and CSV files; no number is
hard-coded in the model code. Thermodynamic properties, heat capacities and reaction equilibrium constants are taken from Cantera
\cite{goodwin2024} with the GRI-Mech~3.0 species data so that kinetics and energy balance are mutually consistent.

\subsection{Reactor tube model}
\label{sec:tube}

The tube is described by a steady one-dimensional pseudo-homogeneous plug-flow model with axial coordinate $z$ from the inlet
(top) to the outlet. The state consists of the molar flows $F_i$ of CH$_4$, H$_2$O, H$_2$, CO, CO$_2$ and N$_2$, the process-gas
temperature $T$ and the pressure $p$. The intrinsic kinetics are those of Xu and Froment \cite{xu1989a} for steam reforming (r1),
water-gas shift (r2) and overall reforming to CO$_2$ (r3):
\begin{align}
r_1 &= \frac{k_1}{p_{\mathrm{H_2}}^{2.5}}\,\frac{p_{\mathrm{CH_4}}\,p_{\mathrm{H_2O}} - p_{\mathrm{H_2}}^{3}\,p_{\mathrm{CO}}/K_1}{\mathrm{DEN}^2}, \\
r_2 &= \frac{k_2}{p_{\mathrm{H_2}}}\,\frac{p_{\mathrm{CO}}\,p_{\mathrm{H_2O}} - p_{\mathrm{H_2}}\,p_{\mathrm{CO_2}}/K_2}{\mathrm{DEN}^2}, \\
r_3 &= \frac{k_3}{p_{\mathrm{H_2}}^{3.5}}\,\frac{p_{\mathrm{CH_4}}\,p_{\mathrm{H_2O}}^{2} - p_{\mathrm{H_2}}^{4}\,p_{\mathrm{CO_2}}/K_3}{\mathrm{DEN}^2}, \\
\mathrm{DEN} &= 1 + K_{\mathrm{CO}}\,p_{\mathrm{CO}} + K_{\mathrm{H_2}}\,p_{\mathrm{H_2}} + K_{\mathrm{CH_4}}\,p_{\mathrm{CH_4}} + K_{\mathrm{H_2O}}\,p_{\mathrm{H_2O}}/p_{\mathrm{H_2}},
\end{align}
with partial pressures in bar. The rate coefficients $k_j$ and adsorption constants $K_i$ follow the Arrhenius and van't Hoff forms
with the pre-exponential factors, activation energies and adsorption enthalpies of Table~5 of \cite{xu1989a}, re-parameterised at
the reference temperatures 648 K and 823 K used in that table so that the reported confidence intervals can be propagated directly
(Section~\ref{sec:uq}). The equilibrium constants $K_1$--$K_3$ are computed from Gibbs energies of formation rather than from the
fitted expressions of the original paper, which removes a source of inconsistency with the energy balance. A floor of
\SI{e-6}{bar} on $p_{\mathrm{H_2}}$ avoids the singularity of the rate expressions at a hydrogen-free inlet.

The balances per tube are
\begin{align}
\frac{\mathrm{d}F_i}{\mathrm{d}z} &= A_{cs}\,\rho_b \sum_j \nu_{ij}\,\eta_j\,r_j, \label{eq:species}\\
\sum_i F_i\,c_{p,i}\,\frac{\mathrm{d}T}{\mathrm{d}z} &= \pi d_i\,U\,\bigl(T_{wo}(z) - T\bigr) - A_{cs}\,\rho_b \sum_j \eta_j\,r_j\,\Delta H_j(T), \label{eq:energy}\\
-\frac{\mathrm{d}p}{\mathrm{d}z} &= \frac{1-\phi}{\phi^{3}}\left(1.75 + 150\,\frac{1-\phi}{Re_p}\right)\frac{\rho_g v_s^{2}}{d_p}, \label{eq:ergun}
\end{align}
where $A_{cs}$ is the cross-section, $\rho_b$ the bed density, $\nu_{ij}$ the stoichiometric coefficients, $\eta_j$ the effectiveness
factors, $\phi$ the bed voidage, $d_p$ the equivalent particle diameter, $v_s$ the superficial velocity and $Re_p$ the particle
Reynolds number; Eq.~\eqref{eq:ergun} is the Ergun equation \cite{ergun1952}. Intraparticle diffusion limitation is severe in
reformer catalysts, and rather than solving the pellet problem we use the constant effectiveness factors adopted by Latham et al.
\cite{latham2011}: $\eta_j = 0.1$ along the tube and 0.05 in the first \SI{10}{\percent} of the length, where the feed contains
no hydrogen and the rate expressions are stiff \todo{confirm the 0.05/0.1 profile and its justification against Latham et al. (2011)}.
Higher alkanes in the feed (C$_2$--C$_5$) are converted to CO and H$_2$ at the inlet by the rule used by Latham
\cite{latham2008} \todo{state the exact inlet conversion rule and confirm its stoichiometry}.

The overall heat-transfer coefficient referred to the inner surface combines the bed-side film coefficient $\alpha_i$ and the
conduction through the tube wall,
\begin{equation}
\frac{1}{U} = \frac{1}{\alpha_i} + \frac{d_i}{2\lambda_t}\ln\frac{d_o}{d_i}, \qquad
T_{wi} = T_{wo} - q''_o\,\frac{d_o}{2\lambda_t}\ln\frac{d_o}{d_i},
\label{eq:U}
\end{equation}
with tube conductivity $\lambda_t = \SI{29.6}{W.m^{-1}.K^{-1}}$ from \cite{latham2008}. The film coefficient uses the Leva--Grummer
correlation \cite{leva1948} in the form written by Latham et al. \cite{latham2011}, multiplied by a heat-transfer factor
$f_{htg}$ which is estimated from plant data (Section~\ref{sec:calibration}). The alternative chain of Xu and Froment
\cite{xu1989b}, based on De~Wasch and Froment \cite{dewasch1972} and Kunii and Smith \cite{kunii1960}, is implemented as an option
and discussed in Section~\ref{sec:verification}. Equations~\eqref{eq:species}--\eqref{eq:ergun} are integrated with the LSODA solver of
SciPy for a prescribed outer-wall profile $T_{wo}(z)$ (stand-alone mode) or as part of the coupled system below.

\subsection{Furnace model and coupling}
\label{sec:furnace}

The furnace is a top-fired box with rows of burners in the roof between rows of tubes (Fig.~\ref{fig:schematic}a; Latham's geometry
has 7 rows of 48 tubes and 8 rows of 12 burners \cite{latham2008}). Instead of the zone method of Hottel and Sarofim
\cite{hottel1967} used by Plehiers and Froment \cite{plehiers1989}, we use a one-dimensional ``long-furnace'' model in the sense of
Roesler \cite{roesler1967}: the flue gas is a plug flow travelling downward co-currently with the process gas, and the radiative
exchange between the gas and the tube surface is lumped into one effective exchange factor $F_{gt}$. Per unit height $z$ (zero at
the roof),
\begin{align}
\dot n_{fg}\,c_{p,fg}\,\frac{\mathrm{d}T_{fg}}{\mathrm{d}z} &= r(z) - q(z)\,\pi d_o N_t, \label{eq:fg}\\
q(z) &= \sigma F_{gt}\bigl(T_{fg}^{4} - T_{wo}^{4}\bigr) + h_c\bigl(T_{fg} - T_{wo}\bigr), \label{eq:q}\\
q(z) &= U\,\bigl(T_{wo} - T\bigr)\,\frac{d_i}{d_o}, \label{eq:coupling}
\end{align}
where $N_t$ is the number of tubes, $q$ the flux per unit outer tube area and $h_c$ a convective coefficient from a flat-plate
correlation \todo{state the convective correlation and its magnitude relative to radiation}. The heat-release density $r(z)$ is a
downward-opening parabola that is zero at $z = L_q L$, releases the fraction $\alpha_{top}$ of the heat in the top fifteenth of the
furnace height and integrates to $(1-f_{loss})\,Q_{comb}$, with $f_{loss} = 0.02$ for casing losses; this is the continuous form of
the sectional heat-release profile of Latham \cite{latham2008}. The dimensionless flame length $L_q$ and $\alpha_{top}$ are
calibration parameters. Because a pure plug flow starts at the adiabatic flame temperature of the fuel mixture and would give an
unphysical cold inlet region before the first heat release, the top zone is treated as well mixed: the fraction $\alpha_{top}$ of
the heat is released instantly at the roof. The combustion heat $Q_{comb}$ and the flue-gas flow and composition follow from the
lower heating value and complete combustion of the fuel gas and pressure-swing-adsorption off-gas with the specified excess air.

Equation~\eqref{eq:coupling} closes the system: at every $z$ the outer-wall temperature $T_{wo}$ is the root of the flux balance
between Eqs.~\eqref{eq:q} and \eqref{eq:coupling}, found by a bracketed root search inside the right-hand side of the ODE system, so
that the furnace and tube equations are integrated together in one downward sweep. The global energy closure of the coupled solution
is verified to $10^{-5}$ of the combustion heat.

\subsection{Tube wall and creep life}
\label{sec:creep}

The creep-life module converts the outer-wall temperature profile into a life estimate. The hoop stress is computed from the
mean-diameter formula of API 530 \cite{api530}, $\sigma_h = p\,(d_o - t)/(2t)$ with wall thickness $t$, evaluated at the pressure of
the process gas at each elevation; thin-wall and Lam\'e variants are implemented for comparison and differ by less than
\todo{quote the difference} at the tube geometry used. The wall thickness is not given by Latham and is taken as \SI{15}{mm}
\todo{confirm a typical wall thickness for a 12.5 m HP-Nb tube of this diameter} and varied in the uncertainty analysis. The
temperature entering the life calculation is the outer-wall temperature, which is conservative with respect to the mid-wall
temperature; the difference is shown in Fig.~\ref{fig:creepprofile}.

The rupture time follows the Larson--Miller parameter \cite{larson1952},
\begin{equation}
\mathrm{LMP} = 10^{-3}\,T\,\bigl(C + \log_{10} t_r\bigr), \qquad \log_{10}\sigma = P_3(\mathrm{LMP}),
\label{eq:lmp}
\end{equation}
with $T$ in K, $t_r$ in h, $C = 22.96$ and a cubic master curve $P_3$ fitted to the minimum-property rupture points of the
manufacturer's curve for Manaurite XM reproduced by Yeh \cite{yeh2021}. This curve is a \emph{placeholder}: the alloy of the
Latham plant is not stated, and the manufacturer's curve carries no scatter information. The module is written so that the curve
can be replaced by a heat-normalised fit to the NIMS creep data sheets for the 25Cr--35Ni--Nb alloys \cite{nims1990,nims1991}
(Section~\ref{sec:data}). Damage under a piecewise-constant history is accumulated with the linear life-fraction rule of Robinson
\cite{robinson1952}, $D = \sum_k \Delta t_k / t_r(T_k, \sigma_k)$, and rupture is predicted at $D = 1$. Throughout the paper, life is
reported as a \emph{consumption rate} $1/t_r$ at the hot spot, or as damage per kmol of hydrogen, relative to the calibrated base
case, because ratios of rupture times are far less sensitive to the placeholder curve than absolute values are.

\subsection{Data}
\label{sec:data}

\paragraph{Verification data} The industrial simulation of Xu and Froment \cite{xu1989b} (Fig.~3 of that paper) provides axial
profiles of methane and CO$_2$ conversion, process-gas, inner-wall and outer-wall temperatures and pressure for one operating point
with the feed of their Table~2. The curves were digitised from the published figure \todo{name the digitising tool and the
resolution} and stored as tidy CSV files with the digitisation metadata (Supplementary Section~S1). The outer-wall curve serves as
the boundary condition and the remaining curves as targets.

\paragraph{Plant data} Latham \cite{latham2008} reports, in Appendix H, four operating points of an industrial top-fired reformer
(Plants A, B and two states of Plant C) with inlet conditions, feed and fuel compositions, outlet temperature and composition,
flue-gas outlet temperature and the tube-wall temperatures measured by infrared pyrometer at two peep-hole elevations (3.66 and
\SI{8.53}{m} from the top of the \SI{12.5}{m} heated length). Table~\ref{tab:latham} lists the cases. The stated measurement
uncertainties are \SI{2}{K} for the outlet temperature, \SI{3}{K} for the wall temperatures and \SI{8}{K} for the flue gas
\todo{confirm these uncertainties against the thesis and whether they are one-sigma}. Plant C1 is excluded from calibration because
its feed flow is inconsistent (equal to the air flow and \SI{22}{\percent} larger than Plant A at a lower rate) and is likely a
transcription error in the source \todo{confirm the exclusion of Plant C1}. The excess air of each case is derived from the
reported fuel, off-gas and air flows and is \SI{21.6}{\percent} for Plant A.

\paragraph{Creep data} The placeholder Larson--Miller curve is transcribed from Yeh \cite{yeh2021}. A transcription template and a
schema for the NIMS data sheets are provided, and an ingestion routine fits $C$ and the master curve by cross-validation with
per-heat offsets; it is validated on a synthetic data set with known $C$ and scatter (recovered within \SI{10}{\percent}) but has not
yet been applied to real NIMS data, which must be transcribed from the printed sheets \todo{state the status of the NIMS
transcription at submission}.

\subsection{Calibration and cross-validation}
\label{sec:calibration}

Four parameters are uncertain in the coupled model and cannot be taken from the literature: the exchange factor $F_{gt}$, the flame
length $L_q$, the top-zone release fraction $\alpha_{top}$ and the heat-transfer factor $f_{htg}$. They are estimated by weighted
non-linear least squares on the three plant cases, with residuals in the outlet temperature, the outlet pressure, the outlet wet
composition, the flue-gas outlet temperature and the two wall temperatures (nine residuals per case), each scaled by its
measurement uncertainty. The trust-region-reflective algorithm with bounds is started from a Latin-hypercube set of initial points
to guard against local minima, and the parameter covariance is obtained from the Jacobian at the optimum with a $t$-based
\SI{95}{\percent} interval. Predictive performance is assessed by leave-one-case-out cross-validation: the parameters are refitted
on two cases and the held-out case is predicted. The four-parameter fit places $\alpha_{top}$ at its bound with a correlation of
0.985 to $L_q$, so the adopted configuration fixes $\alpha_{top} = 0.182$ \todo{state the origin of the value 0.182} and estimates
the remaining three; the increase in $\chi^2$ is reported in Section~\ref{sec:verification}.

\subsection{Surrogates and conformal prediction}
\label{sec:surrogate}

The coupled model takes 1--3 s per operating point, which is adequate for calibration but not for hourly annual histories,
evolutionary optimisation or Monte Carlo. Surrogates are trained on a Latin-hypercube design of \nLHS{} points over the seven
operating inputs (Table~S1 of the Supplementary material) \todo{decide whether the ranges table goes in the main text or the supplement} for eight
outputs (hot-spot temperature $T_{wo,max}$ and its position, outlet temperature and pressure, dry methane slip, net hydrogen, duty,
and $\log_{10} t_r$ at the hot spot). Three model families are compared with scikit-learn \cite{pedregosa2011}: Gaussian-process
regression with a Mat\'ern 5/2 kernel on a 1000-point subset, histogram gradient boosting and a multilayer perceptron, each with and
without two physics-derived features (the equilibrium methane conversion at outlet conditions and the specific duty). The model
with the lowest five-fold cross-validated error is selected per target and evaluated on an independent Saltelli design of
\nSaltelli{} points. Prediction intervals are obtained by the cross-conformal (CV+) method \cite{barber2021} as implemented in MAPIE
\cite{taquet2022}, which provides finite-sample marginal coverage without a distributional assumption; the Gaussian-process
native intervals are reported for comparison.

\subsection{Sensitivity, scenarios, optimisation and uncertainty propagation}
\label{sec:campaigns}

\paragraph{RQ1} Global sensitivity uses the variance-based Sobol indices \cite{saltelli2010} estimated with SALib
\cite{herman2017} on the Saltelli design evaluated with the physics model, together with an open-loop Latin-hypercube campaign and a
closed-loop grid in steam-to-carbon ratio and load in which the firing is solved so that the outlet temperature equals the Plant A
value (\SI{1105.55}{K}).

\paragraph{RQ2} Annual histories of 8760 hourly quasi-steady states are evaluated through the surrogates, with every 24th hour and
every hour outside the training bounds re-evaluated by the physics model. Two control modes are represented by inverting the
surrogates for the firing factor with a vectorised bisection: holding the outlet temperature at its base value, or holding the dry
methane slip. The scenarios are S1 steady base; S2 daily load-following (full load 06--22 h, \SI{70}{\percent} at night, ramps of
0.10 per hour); S3 renewable-following (a bounded Ornstein--Uhlenbeck load process between 0.60 and 1.10 with mean 0.85); S4
over-firing events (24 events per year of \SI{+10}{\percent} firing for 2 h, six events of \SI{+25}{\percent} for 1 h, and a
flame-impingement proxy adding \SI{40}{K} to the hot spot for 24 h per year); S5 catalyst ageing, with the activity falling
linearly from 0.30 to 0.10 over four years under each control mode; and S6 steady operation at steam-to-carbon 2.5 and 3.5. Damage
is accumulated hourly by Eq.~\eqref{eq:lmp} and Robinson's rule, and each scenario is also evaluated at its annual-mean inputs to
quantify the error of a steady-state assessment. Start-up and shutdown transients, thermal fatigue and creep--fatigue interaction are
outside the scope.

\paragraph{RQ3} The non-dominated sorting genetic algorithm NSGA-II \cite{deb2002} in pymoo \cite{blank2020} maximises net hydrogen
per tube and minimises firebox fuel per kmol H$_2$ and life consumption per kmol H$_2$ over five decision variables (load,
steam-to-carbon ratio, inlet temperature, excess air, firing) subject to methane slip not above the base value, hot-spot
temperature not above \SI{1193}{K} \todo{justify the 1193 K limit, e.g. the design metal temperature of the alloy} and outlet
temperature not above \SI{1125}{K}, using the Gaussian-process surrogates (population 200, 300 generations). Every Pareto member is
re-evaluated with the physics model and the constraints are re-checked. Because firebox fuel per kmol H$_2$ ignores the cost of
raising the additional steam, a \emph{steam credit} $\alpha$ is introduced: the total fuel counts firebox fuel plus a fraction
$(1-\alpha)$ of the heat needed to raise the process steam and preheat the feed, so that $\alpha = 1$ recovers the firebox-only
objective and $\alpha = 0$ charges steam in full \todo{confirm the steam-credit definition and the enthalpy basis for steam raising}.
The iso-production front (hydrogen within \SI{1}{\percent} of the base) is recomputed for $\alpha$ from 0 to 1.

\paragraph{RQ4} \label{sec:uq} Uncertain inputs are grouped into kinetics (the six rate parameters and four adsorption constants of
Xu and Froment with normal distributions from their reported \SI{95}{\percent} intervals, a log-uniform factor of 0.5--2 on the
effectiveness factors and $\pm 0.03$ on the catalyst activity), heat transfer (the calibrated $F_{gt}$, $L_q$ and $f_{htg}$ with
their joint covariance), creep (a normal scatter of 0.3 decades on $\log_{10} t_r$ representing a factor of two between heats, $C$
uniform within $\pm 1$ of 22.96 with the master curve refitted for each $C$, wall thickness 12--18 mm and tube conductivity
$\pm 15$\,\%) and measurement (a uniform $\pm 10$ K offset on the pyrometer-referenced hot-spot temperature). A scrambled Sobol
design of \nMC{} samples is propagated through the physics model at eight operating points taken from RQ2 and RQ3; the same sample
is used at every point so that ratios between points are paired. Variance is decomposed by varying one group at a time (1024
samples per group) and, independently, by Sobol indices computed on a Gaussian-process surrogate of the joint sample. Finally, the
fraction of paired samples for which each qualitative conclusion of RQ2 and RQ3 holds is reported, both with the calibrated flame
length and with a structural variant in which the flame length shrinks with load as $L_q\,\mathrm{load}^{0.5}$
\todo{justify or replace the load exponent 0.5 of the flame-length variant}. The complete pipeline from design to figures runs in
\pipelineMin{} minutes on a laptop.
